% ============================================================
%  课程笔记封面模板 —— 「千里江山图」风格
%  取材：王希孟《千里江山图》（北宋·青绿山水长卷）：赭黄绢底上以
%        石青、石绿层层积染的群峰，留白的江天与细笔水纹。
%  与 cover_Impression风.tex 同构（眉题 / 主标题 / 菱形饰线 /
%  副标题 / 底部作者日期），直接 \input{} 复用：
%      % ============================================================
%  课程笔记封面模板 —— 「千里江山图」风格
%  取材：王希孟《千里江山图》（北宋·青绿山水长卷）：赭黄绢底上以
%        石青、石绿层层积染的群峰，留白的江天与细笔水纹。
%  与 cover_Impression风.tex 同构（眉题 / 主标题 / 菱形饰线 /
%  副标题 / 底部作者日期），直接 \input{} 复用：
%      % ============================================================
%  课程笔记封面模板 —— 「千里江山图」风格
%  取材：王希孟《千里江山图》（北宋·青绿山水长卷）：赭黄绢底上以
%        石青、石绿层层积染的群峰，留白的江天与细笔水纹。
%  与 cover_Impression风.tex 同构（眉题 / 主标题 / 菱形饰线 /
%  副标题 / 底部作者日期），直接 \input{} 复用：
%      % ============================================================
%  课程笔记封面模板 —— 「千里江山图」风格
%  取材：王希孟《千里江山图》（北宋·青绿山水长卷）：赭黄绢底上以
%        石青、石绿层层积染的群峰，留白的江天与细笔水纹。
%  与 cover_Impression风.tex 同构（眉题 / 主标题 / 菱形饰线 /
%  副标题 / 底部作者日期），直接 \input{} 复用：
%      \input{cover_千里江山图风}
%
%  配色：ql* 前缀（青绿设色）。本文件用 \providecolor 兜底，
%        单独复制到任何笔记本也能编译；想整体换调改这里即可。
%
%  注意：整幅画面绘于 \begin{scope}[shift={(current page.south west)}] 内，
%        所有坐标以「页面左下角为原点、单位 cm」书写；脱离该 scope
%        直接写 (x,y)，图形会落到页面外而不可见。
%
%  可用字段（在主文件导言区用 \newcommand 预定义即可覆盖缺省值）：
%      \notename     主标题（必填，主模板已定义）
%      \noteauthor   作者（必填，主模板已定义）
%      \notecourseen 眉题英文（缺省 Course Notes）
%      \notesubtitle 副标题（缺省「学习笔记」）
%      \noteextra    底部附加信息（缺省不显示）
% ============================================================

% ---- 《千里江山图》取色（青绿设色） ----
\providecolor{qlSilk}{RGB}{230,218,184}     % 绢底：赭黄
\providecolor{qlSilkDark}{RGB}{208,192,152} % 绢底：暗部
\providecolor{qlGreen2}{RGB}{150,186,128}   % 石绿（淡，远峰）
\providecolor{qlGreen}{RGB}{82,140,106}     % 石绿（主峰）
\providecolor{qlBlue}{RGB}{40,94,132}       % 石青（主峰阴面）
\providecolor{qlOchre}{RGB}{178,134,86}     % 赭石（坡脚/滩渚）
\providecolor{qlInk}{RGB}{56,54,50}         % 墨（点缀/标题）
\providecolor{qlSun}{RGB}{238,214,150}      % 微茫的日色

% ---- 缺省字段（主文件已定义的同名命令不会被覆盖） ----
\providecommand{\notecourseen}{Course Notes}
\providecommand{\notesubtitle}{学\ 习\ 笔\ 记}
\providecommand{\noteextra}{}

\begin{titlepage}
	\begin{tikzpicture}[remember picture, overlay]
		\begin{scope}[shift={(current page.south west)}]

		% ============================================================
		%  1. 绢底（赭黄，下缘略深）
		% ============================================================
		\shade[top color=qlSilk, bottom color=qlSilkDark] (0,0) rectangle (21,29.7);

		% ============================================================
		%  2. 微茫的日色
		% ============================================================
		\fill[qlSun, opacity=0.35] (16.6,26.2) circle[radius=2.6];
		\fill[qlSun, opacity=0.18] (16.6,26.2) circle[radius=4.0];

		% ============================================================
		%  3. 远峰（淡石绿，没入江天）
		% ============================================================
		\fill[qlGreen2, opacity=0.45]
			(0,15.0) -- (2.2,16.2) -- (3.8,15.2) -- (5.6,16.4) -- (7.6,15.0)
			-- (9.8,16.6) -- (12.2,15.2) -- (14.6,16.4) -- (17.0,15.0)
			-- (19.2,16.0) -- (21,14.8) -- (21,12.4) -- (0,12.4) -- cycle;

		% ============================================================
		%  4. 主峰群（石绿 → 石青，自远及近）
		% ============================================================
		% 左峰
		\shade[top color=qlGreen2, bottom color=qlGreen]
			(0,8.4) -- (1.4,10.6) -- (2.6,12.6) -- (3.6,13.6) -- (4.6,12.0)
			-- (5.6,13.0) -- (6.8,11.2) -- (8.0,9.4) -- (9.0,8.2)
			-- (9.4,6.4) -- (9.0,0) -- (0,0) -- cycle;
		% 中峰（最高）
		\shade[top color=qlGreen, bottom color=qlBlue]
			(7.6,8.6) -- (8.8,11.4) -- (9.8,13.8) -- (10.6,15.4) -- (11.4,16.2)
			-- (12.4,14.6) -- (13.4,15.6) -- (14.4,13.2) -- (15.4,11.0)
			-- (16.2,9.0) -- (16.6,7.0) -- (16.4,0) -- (7.6,0) -- cycle;
		% 右峰
		\shade[top color=qlGreen2, bottom color=qlGreen]
			(15.8,8.6) -- (16.8,10.4) -- (17.8,12.2) -- (18.6,13.4) -- (19.6,11.8)
			-- (20.4,12.6) -- (21,10.8) -- (21,0) -- (15.8,0) -- cycle;

		% ============================================================
		%  5. 山体皴线与峰顶苔点
		% ============================================================
		\foreach \x/\y/\l in {2.2/10.4/1.6, 3.6/12.2/1.8, 5.2/11.6/1.4,
		                      9.6/12.6/2.0, 11.0/14.6/2.2, 12.8/13.8/1.8,
		                      14.6/11.6/1.6, 17.6/11.0/1.5, 19.2/11.4/1.3}{
			\draw[qlBlue, opacity=0.30, line width=0.9pt]
				(\x,\y) .. controls (\x+0.22,\y-\l*0.5) and (\x-0.14,\y-\l*0.8) .. (\x+0.12,\y-\l);
		}
		\foreach \x/\y in {3.6/13.7, 5.6/13.1, 9.8/13.9, 11.4/16.3, 13.4/15.7,
		                   18.6/13.5, 20.4/12.7}{
			\fill[qlInk, opacity=0.55] (\x,\y) ellipse[x radius=0.22, y radius=0.12];
		}

		% ============================================================
		%  6. 江水（留白 + 细笔水纹）
		% ============================================================
		\foreach \y in {5.8, 5.0, 4.2, 3.4}{
			\draw[qlBlue, opacity=0.28, line width=0.7pt]
				plot[domain=0.8:20.2, samples=44, smooth]
					(\x, {\y + 0.11*sin(deg(2.4*\x + 4*\y))});
		}

		% ============================================================
		%  7. 滩渚、村落与渔舟
		% ============================================================
		\fill[qlOchre, opacity=0.75]
			(0,0) -- (0,2.4) -- (2.0,3.0) -- (4.2,2.2) -- (6.4,2.8)
			-- (8.6,2.0) -- (10.8,2.6) -- (13.0,1.8) -- (15.4,2.4)
			-- (17.8,1.6) -- (21,2.2) -- (21,0) -- cycle;
		\foreach \x/\y in {2.4/3.0, 3.1/3.0, 2.75/3.5, 13.6/2.5, 14.3/2.5, 13.95/3.0}{
			\draw[qlInk, opacity=0.6, line width=0.8pt] (\x,\y) -- (\x,\y+0.62);
			\fill[qlInk, opacity=0.5] (\x,\y+0.72) circle[radius=0.15];
		}
		\fill[qlInk, opacity=0.85]
			(10.4,4.5) .. controls (11.0,4.28) and (11.9,4.28) .. (12.5,4.5)
			.. controls (12.0,4.74) and (10.9,4.74) .. (10.4,4.5) -- cycle;
		\draw[qlInk, opacity=0.85, line width=0.8pt] (11.45,4.5) -- (11.45,5.2);

		% ============================================================
		%  8. 标题衬底（一层薄绢，保证文字可读）
		% ============================================================
		\fill[qlSilk, opacity=0.55]
			(2.6,16.9) -- (17.4,17.3) -- (18.0,22.6) -- (16.4,23.6) -- (2.2,22.9) -- cycle;

		% ============================================================
		%  9. 细边框（墨线）
		% ============================================================
		\draw[qlInk, line width=0.7pt, opacity=0.45]
			([shift={(0.85,0.85)}]current page.south west) rectangle
			([shift={(-0.85,-0.85)}]current page.north east);

		% ============================================================
		% 10. 顶部眉题（课程英文小字 + 自适应墨线）
		% ============================================================
		\node[anchor=north, text=qlInk, align=center]
			(qlEyebrow) at ([yshift=-3.4cm]current page.north) {%
			{\sffamily\fontsize{12}{15}\selectfont \uppercase{\notecourseen}}
		};
		\draw[qlInk, line width=1.0pt] ($(qlEyebrow.west)+(-0.25cm,0)$) -- ++(-1.1cm,0);
		\draw[qlInk, line width=1.0pt] ($(qlEyebrow.east)+(0.25cm,0)$) -- ++(1.1cm,0);

		% ============================================================
		% 11. 主标题（居中，使用 \notename）
		% ============================================================
		\node[anchor=center, align=center, text=qlInk]
			at ([yshift=-8.3cm]current page.north) {%
			{\fontsize{38}{48}\sffamily\bfseries \notename}
		};

		% ============================================================
		% 12. 菱形 + 双细线
		% ============================================================
		\coordinate (C) at ([yshift=-11.6cm]current page.north);
		\draw[qlInk, line width=1.1pt] (C) ++(-2.2cm,0) -- ++(1.55cm,0);
		\draw[qlInk, line width=1.1pt] (C) ++(0.65cm,0) -- ++(1.55cm,0);
		\node[fill=qlInk, minimum size=0.24cm, inner sep=0pt, rotate=45] at (C) {};

		% ============================================================
		% 13. 副标题（使用 \notesubtitle）
		% ============================================================
		\node[anchor=north, text=qlInk, align=center]
			at ([yshift=-12.4cm]current page.north) {
			{\fontsize{15}{20}\sffamily \notesubtitle}
		};

		% ============================================================
		% 14. 底部作者、日期、附加信息
		% ============================================================
		\coordinate (B) at ([yshift=2.7cm]current page.south);
		\draw[qlInk, line width=0.8pt, opacity=0.8] (B) ++(-1.1cm,0.95cm) -- ++(2.2cm,0);
		\node[anchor=south, text=qlInk, align=center] at (B) {
			\begin{tabular}{c}
				{\fontsize{19}{24}\sffamily\bfseries \noteauthor} \\[0.35cm]
				{\large \sffamily \today}
				\ifstrempty{\noteextra}{}{\\[0.3cm]{\normalsize\sffamily \noteextra}}
			\end{tabular}
		};

		% ============================================================
		% 15. 右下角落款（画作信息）
		% ============================================================
		\node[anchor=south east, text=qlInk, opacity=0.62, font=\footnotesize\itshape]
			at ([shift={(-1.35,1.25)}]current page.south east)
			{Wang Ximeng, A Thousand Li of Rivers and Mountains, 1113};

		\end{scope}
	\end{tikzpicture}
\end{titlepage}

%
%  配色：ql* 前缀（青绿设色）。本文件用 \providecolor 兜底，
%        单独复制到任何笔记本也能编译；想整体换调改这里即可。
%
%  注意：整幅画面绘于 \begin{scope}[shift={(current page.south west)}] 内，
%        所有坐标以「页面左下角为原点、单位 cm」书写；脱离该 scope
%        直接写 (x,y)，图形会落到页面外而不可见。
%
%  可用字段（在主文件导言区用 \newcommand 预定义即可覆盖缺省值）：
%      \notename     主标题（必填，主模板已定义）
%      \noteauthor   作者（必填，主模板已定义）
%      \notecourseen 眉题英文（缺省 Course Notes）
%      \notesubtitle 副标题（缺省「学习笔记」）
%      \noteextra    底部附加信息（缺省不显示）
% ============================================================

% ---- 《千里江山图》取色（青绿设色） ----
\providecolor{qlSilk}{RGB}{230,218,184}     % 绢底：赭黄
\providecolor{qlSilkDark}{RGB}{208,192,152} % 绢底：暗部
\providecolor{qlGreen2}{RGB}{150,186,128}   % 石绿（淡，远峰）
\providecolor{qlGreen}{RGB}{82,140,106}     % 石绿（主峰）
\providecolor{qlBlue}{RGB}{40,94,132}       % 石青（主峰阴面）
\providecolor{qlOchre}{RGB}{178,134,86}     % 赭石（坡脚/滩渚）
\providecolor{qlInk}{RGB}{56,54,50}         % 墨（点缀/标题）
\providecolor{qlSun}{RGB}{238,214,150}      % 微茫的日色

% ---- 缺省字段（主文件已定义的同名命令不会被覆盖） ----
\providecommand{\notecourseen}{Course Notes}
\providecommand{\notesubtitle}{学\ 习\ 笔\ 记}
\providecommand{\noteextra}{}

\begin{titlepage}
	\begin{tikzpicture}[remember picture, overlay]
		\begin{scope}[shift={(current page.south west)}]

		% ============================================================
		%  1. 绢底（赭黄，下缘略深）
		% ============================================================
		\shade[top color=qlSilk, bottom color=qlSilkDark] (0,0) rectangle (21,29.7);

		% ============================================================
		%  2. 微茫的日色
		% ============================================================
		\fill[qlSun, opacity=0.35] (16.6,26.2) circle[radius=2.6];
		\fill[qlSun, opacity=0.18] (16.6,26.2) circle[radius=4.0];

		% ============================================================
		%  3. 远峰（淡石绿，没入江天）
		% ============================================================
		\fill[qlGreen2, opacity=0.45]
			(0,15.0) -- (2.2,16.2) -- (3.8,15.2) -- (5.6,16.4) -- (7.6,15.0)
			-- (9.8,16.6) -- (12.2,15.2) -- (14.6,16.4) -- (17.0,15.0)
			-- (19.2,16.0) -- (21,14.8) -- (21,12.4) -- (0,12.4) -- cycle;

		% ============================================================
		%  4. 主峰群（石绿 → 石青，自远及近）
		% ============================================================
		% 左峰
		\shade[top color=qlGreen2, bottom color=qlGreen]
			(0,8.4) -- (1.4,10.6) -- (2.6,12.6) -- (3.6,13.6) -- (4.6,12.0)
			-- (5.6,13.0) -- (6.8,11.2) -- (8.0,9.4) -- (9.0,8.2)
			-- (9.4,6.4) -- (9.0,0) -- (0,0) -- cycle;
		% 中峰（最高）
		\shade[top color=qlGreen, bottom color=qlBlue]
			(7.6,8.6) -- (8.8,11.4) -- (9.8,13.8) -- (10.6,15.4) -- (11.4,16.2)
			-- (12.4,14.6) -- (13.4,15.6) -- (14.4,13.2) -- (15.4,11.0)
			-- (16.2,9.0) -- (16.6,7.0) -- (16.4,0) -- (7.6,0) -- cycle;
		% 右峰
		\shade[top color=qlGreen2, bottom color=qlGreen]
			(15.8,8.6) -- (16.8,10.4) -- (17.8,12.2) -- (18.6,13.4) -- (19.6,11.8)
			-- (20.4,12.6) -- (21,10.8) -- (21,0) -- (15.8,0) -- cycle;

		% ============================================================
		%  5. 山体皴线与峰顶苔点
		% ============================================================
		\foreach \x/\y/\l in {2.2/10.4/1.6, 3.6/12.2/1.8, 5.2/11.6/1.4,
		                      9.6/12.6/2.0, 11.0/14.6/2.2, 12.8/13.8/1.8,
		                      14.6/11.6/1.6, 17.6/11.0/1.5, 19.2/11.4/1.3}{
			\draw[qlBlue, opacity=0.30, line width=0.9pt]
				(\x,\y) .. controls (\x+0.22,\y-\l*0.5) and (\x-0.14,\y-\l*0.8) .. (\x+0.12,\y-\l);
		}
		\foreach \x/\y in {3.6/13.7, 5.6/13.1, 9.8/13.9, 11.4/16.3, 13.4/15.7,
		                   18.6/13.5, 20.4/12.7}{
			\fill[qlInk, opacity=0.55] (\x,\y) ellipse[x radius=0.22, y radius=0.12];
		}

		% ============================================================
		%  6. 江水（留白 + 细笔水纹）
		% ============================================================
		\foreach \y in {5.8, 5.0, 4.2, 3.4}{
			\draw[qlBlue, opacity=0.28, line width=0.7pt]
				plot[domain=0.8:20.2, samples=44, smooth]
					(\x, {\y + 0.11*sin(deg(2.4*\x + 4*\y))});
		}

		% ============================================================
		%  7. 滩渚、村落与渔舟
		% ============================================================
		\fill[qlOchre, opacity=0.75]
			(0,0) -- (0,2.4) -- (2.0,3.0) -- (4.2,2.2) -- (6.4,2.8)
			-- (8.6,2.0) -- (10.8,2.6) -- (13.0,1.8) -- (15.4,2.4)
			-- (17.8,1.6) -- (21,2.2) -- (21,0) -- cycle;
		\foreach \x/\y in {2.4/3.0, 3.1/3.0, 2.75/3.5, 13.6/2.5, 14.3/2.5, 13.95/3.0}{
			\draw[qlInk, opacity=0.6, line width=0.8pt] (\x,\y) -- (\x,\y+0.62);
			\fill[qlInk, opacity=0.5] (\x,\y+0.72) circle[radius=0.15];
		}
		\fill[qlInk, opacity=0.85]
			(10.4,4.5) .. controls (11.0,4.28) and (11.9,4.28) .. (12.5,4.5)
			.. controls (12.0,4.74) and (10.9,4.74) .. (10.4,4.5) -- cycle;
		\draw[qlInk, opacity=0.85, line width=0.8pt] (11.45,4.5) -- (11.45,5.2);

		% ============================================================
		%  8. 标题衬底（一层薄绢，保证文字可读）
		% ============================================================
		\fill[qlSilk, opacity=0.55]
			(2.6,16.9) -- (17.4,17.3) -- (18.0,22.6) -- (16.4,23.6) -- (2.2,22.9) -- cycle;

		% ============================================================
		%  9. 细边框（墨线）
		% ============================================================
		\draw[qlInk, line width=0.7pt, opacity=0.45]
			([shift={(0.85,0.85)}]current page.south west) rectangle
			([shift={(-0.85,-0.85)}]current page.north east);

		% ============================================================
		% 10. 顶部眉题（课程英文小字 + 自适应墨线）
		% ============================================================
		\node[anchor=north, text=qlInk, align=center]
			(qlEyebrow) at ([yshift=-3.4cm]current page.north) {%
			{\sffamily\fontsize{12}{15}\selectfont \uppercase{\notecourseen}}
		};
		\draw[qlInk, line width=1.0pt] ($(qlEyebrow.west)+(-0.25cm,0)$) -- ++(-1.1cm,0);
		\draw[qlInk, line width=1.0pt] ($(qlEyebrow.east)+(0.25cm,0)$) -- ++(1.1cm,0);

		% ============================================================
		% 11. 主标题（居中，使用 \notename）
		% ============================================================
		\node[anchor=center, align=center, text=qlInk]
			at ([yshift=-8.3cm]current page.north) {%
			{\fontsize{38}{48}\sffamily\bfseries \notename}
		};

		% ============================================================
		% 12. 菱形 + 双细线
		% ============================================================
		\coordinate (C) at ([yshift=-11.6cm]current page.north);
		\draw[qlInk, line width=1.1pt] (C) ++(-2.2cm,0) -- ++(1.55cm,0);
		\draw[qlInk, line width=1.1pt] (C) ++(0.65cm,0) -- ++(1.55cm,0);
		\node[fill=qlInk, minimum size=0.24cm, inner sep=0pt, rotate=45] at (C) {};

		% ============================================================
		% 13. 副标题（使用 \notesubtitle）
		% ============================================================
		\node[anchor=north, text=qlInk, align=center]
			at ([yshift=-12.4cm]current page.north) {
			{\fontsize{15}{20}\sffamily \notesubtitle}
		};

		% ============================================================
		% 14. 底部作者、日期、附加信息
		% ============================================================
		\coordinate (B) at ([yshift=2.7cm]current page.south);
		\draw[qlInk, line width=0.8pt, opacity=0.8] (B) ++(-1.1cm,0.95cm) -- ++(2.2cm,0);
		\node[anchor=south, text=qlInk, align=center] at (B) {
			\begin{tabular}{c}
				{\fontsize{19}{24}\sffamily\bfseries \noteauthor} \\[0.35cm]
				{\large \sffamily \today}
				\ifstrempty{\noteextra}{}{\\[0.3cm]{\normalsize\sffamily \noteextra}}
			\end{tabular}
		};

		% ============================================================
		% 15. 右下角落款（画作信息）
		% ============================================================
		\node[anchor=south east, text=qlInk, opacity=0.62, font=\footnotesize\itshape]
			at ([shift={(-1.35,1.25)}]current page.south east)
			{Wang Ximeng, A Thousand Li of Rivers and Mountains, 1113};

		\end{scope}
	\end{tikzpicture}
\end{titlepage}

%
%  配色：ql* 前缀（青绿设色）。本文件用 \providecolor 兜底，
%        单独复制到任何笔记本也能编译；想整体换调改这里即可。
%
%  注意：整幅画面绘于 \begin{scope}[shift={(current page.south west)}] 内，
%        所有坐标以「页面左下角为原点、单位 cm」书写；脱离该 scope
%        直接写 (x,y)，图形会落到页面外而不可见。
%
%  可用字段（在主文件导言区用 \newcommand 预定义即可覆盖缺省值）：
%      \notename     主标题（必填，主模板已定义）
%      \noteauthor   作者（必填，主模板已定义）
%      \notecourseen 眉题英文（缺省 Course Notes）
%      \notesubtitle 副标题（缺省「学习笔记」）
%      \noteextra    底部附加信息（缺省不显示）
% ============================================================

% ---- 《千里江山图》取色（青绿设色） ----
\providecolor{qlSilk}{RGB}{230,218,184}     % 绢底：赭黄
\providecolor{qlSilkDark}{RGB}{208,192,152} % 绢底：暗部
\providecolor{qlGreen2}{RGB}{150,186,128}   % 石绿（淡，远峰）
\providecolor{qlGreen}{RGB}{82,140,106}     % 石绿（主峰）
\providecolor{qlBlue}{RGB}{40,94,132}       % 石青（主峰阴面）
\providecolor{qlOchre}{RGB}{178,134,86}     % 赭石（坡脚/滩渚）
\providecolor{qlInk}{RGB}{56,54,50}         % 墨（点缀/标题）
\providecolor{qlSun}{RGB}{238,214,150}      % 微茫的日色

% ---- 缺省字段（主文件已定义的同名命令不会被覆盖） ----
\providecommand{\notecourseen}{Course Notes}
\providecommand{\notesubtitle}{学\ 习\ 笔\ 记}
\providecommand{\noteextra}{}

\begin{titlepage}
	\begin{tikzpicture}[remember picture, overlay]
		\begin{scope}[shift={(current page.south west)}]

		% ============================================================
		%  1. 绢底（赭黄，下缘略深）
		% ============================================================
		\shade[top color=qlSilk, bottom color=qlSilkDark] (0,0) rectangle (21,29.7);

		% ============================================================
		%  2. 微茫的日色
		% ============================================================
		\fill[qlSun, opacity=0.35] (16.6,26.2) circle[radius=2.6];
		\fill[qlSun, opacity=0.18] (16.6,26.2) circle[radius=4.0];

		% ============================================================
		%  3. 远峰（淡石绿，没入江天）
		% ============================================================
		\fill[qlGreen2, opacity=0.45]
			(0,15.0) -- (2.2,16.2) -- (3.8,15.2) -- (5.6,16.4) -- (7.6,15.0)
			-- (9.8,16.6) -- (12.2,15.2) -- (14.6,16.4) -- (17.0,15.0)
			-- (19.2,16.0) -- (21,14.8) -- (21,12.4) -- (0,12.4) -- cycle;

		% ============================================================
		%  4. 主峰群（石绿 → 石青，自远及近）
		% ============================================================
		% 左峰
		\shade[top color=qlGreen2, bottom color=qlGreen]
			(0,8.4) -- (1.4,10.6) -- (2.6,12.6) -- (3.6,13.6) -- (4.6,12.0)
			-- (5.6,13.0) -- (6.8,11.2) -- (8.0,9.4) -- (9.0,8.2)
			-- (9.4,6.4) -- (9.0,0) -- (0,0) -- cycle;
		% 中峰（最高）
		\shade[top color=qlGreen, bottom color=qlBlue]
			(7.6,8.6) -- (8.8,11.4) -- (9.8,13.8) -- (10.6,15.4) -- (11.4,16.2)
			-- (12.4,14.6) -- (13.4,15.6) -- (14.4,13.2) -- (15.4,11.0)
			-- (16.2,9.0) -- (16.6,7.0) -- (16.4,0) -- (7.6,0) -- cycle;
		% 右峰
		\shade[top color=qlGreen2, bottom color=qlGreen]
			(15.8,8.6) -- (16.8,10.4) -- (17.8,12.2) -- (18.6,13.4) -- (19.6,11.8)
			-- (20.4,12.6) -- (21,10.8) -- (21,0) -- (15.8,0) -- cycle;

		% ============================================================
		%  5. 山体皴线与峰顶苔点
		% ============================================================
		\foreach \x/\y/\l in {2.2/10.4/1.6, 3.6/12.2/1.8, 5.2/11.6/1.4,
		                      9.6/12.6/2.0, 11.0/14.6/2.2, 12.8/13.8/1.8,
		                      14.6/11.6/1.6, 17.6/11.0/1.5, 19.2/11.4/1.3}{
			\draw[qlBlue, opacity=0.30, line width=0.9pt]
				(\x,\y) .. controls (\x+0.22,\y-\l*0.5) and (\x-0.14,\y-\l*0.8) .. (\x+0.12,\y-\l);
		}
		\foreach \x/\y in {3.6/13.7, 5.6/13.1, 9.8/13.9, 11.4/16.3, 13.4/15.7,
		                   18.6/13.5, 20.4/12.7}{
			\fill[qlInk, opacity=0.55] (\x,\y) ellipse[x radius=0.22, y radius=0.12];
		}

		% ============================================================
		%  6. 江水（留白 + 细笔水纹）
		% ============================================================
		\foreach \y in {5.8, 5.0, 4.2, 3.4}{
			\draw[qlBlue, opacity=0.28, line width=0.7pt]
				plot[domain=0.8:20.2, samples=44, smooth]
					(\x, {\y + 0.11*sin(deg(2.4*\x + 4*\y))});
		}

		% ============================================================
		%  7. 滩渚、村落与渔舟
		% ============================================================
		\fill[qlOchre, opacity=0.75]
			(0,0) -- (0,2.4) -- (2.0,3.0) -- (4.2,2.2) -- (6.4,2.8)
			-- (8.6,2.0) -- (10.8,2.6) -- (13.0,1.8) -- (15.4,2.4)
			-- (17.8,1.6) -- (21,2.2) -- (21,0) -- cycle;
		\foreach \x/\y in {2.4/3.0, 3.1/3.0, 2.75/3.5, 13.6/2.5, 14.3/2.5, 13.95/3.0}{
			\draw[qlInk, opacity=0.6, line width=0.8pt] (\x,\y) -- (\x,\y+0.62);
			\fill[qlInk, opacity=0.5] (\x,\y+0.72) circle[radius=0.15];
		}
		\fill[qlInk, opacity=0.85]
			(10.4,4.5) .. controls (11.0,4.28) and (11.9,4.28) .. (12.5,4.5)
			.. controls (12.0,4.74) and (10.9,4.74) .. (10.4,4.5) -- cycle;
		\draw[qlInk, opacity=0.85, line width=0.8pt] (11.45,4.5) -- (11.45,5.2);

		% ============================================================
		%  8. 标题衬底（一层薄绢，保证文字可读）
		% ============================================================
		\fill[qlSilk, opacity=0.55]
			(2.6,16.9) -- (17.4,17.3) -- (18.0,22.6) -- (16.4,23.6) -- (2.2,22.9) -- cycle;

		% ============================================================
		%  9. 细边框（墨线）
		% ============================================================
		\draw[qlInk, line width=0.7pt, opacity=0.45]
			([shift={(0.85,0.85)}]current page.south west) rectangle
			([shift={(-0.85,-0.85)}]current page.north east);

		% ============================================================
		% 10. 顶部眉题（课程英文小字 + 自适应墨线）
		% ============================================================
		\node[anchor=north, text=qlInk, align=center]
			(qlEyebrow) at ([yshift=-3.4cm]current page.north) {%
			{\sffamily\fontsize{12}{15}\selectfont \uppercase{\notecourseen}}
		};
		\draw[qlInk, line width=1.0pt] ($(qlEyebrow.west)+(-0.25cm,0)$) -- ++(-1.1cm,0);
		\draw[qlInk, line width=1.0pt] ($(qlEyebrow.east)+(0.25cm,0)$) -- ++(1.1cm,0);

		% ============================================================
		% 11. 主标题（居中，使用 \notename）
		% ============================================================
		\node[anchor=center, align=center, text=qlInk]
			at ([yshift=-8.3cm]current page.north) {%
			{\fontsize{38}{48}\sffamily\bfseries \notename}
		};

		% ============================================================
		% 12. 菱形 + 双细线
		% ============================================================
		\coordinate (C) at ([yshift=-11.6cm]current page.north);
		\draw[qlInk, line width=1.1pt] (C) ++(-2.2cm,0) -- ++(1.55cm,0);
		\draw[qlInk, line width=1.1pt] (C) ++(0.65cm,0) -- ++(1.55cm,0);
		\node[fill=qlInk, minimum size=0.24cm, inner sep=0pt, rotate=45] at (C) {};

		% ============================================================
		% 13. 副标题（使用 \notesubtitle）
		% ============================================================
		\node[anchor=north, text=qlInk, align=center]
			at ([yshift=-12.4cm]current page.north) {
			{\fontsize{15}{20}\sffamily \notesubtitle}
		};

		% ============================================================
		% 14. 底部作者、日期、附加信息
		% ============================================================
		\coordinate (B) at ([yshift=2.7cm]current page.south);
		\draw[qlInk, line width=0.8pt, opacity=0.8] (B) ++(-1.1cm,0.95cm) -- ++(2.2cm,0);
		\node[anchor=south, text=qlInk, align=center] at (B) {
			\begin{tabular}{c}
				{\fontsize{19}{24}\sffamily\bfseries \noteauthor} \\[0.35cm]
				{\large \sffamily \today}
				\ifstrempty{\noteextra}{}{\\[0.3cm]{\normalsize\sffamily \noteextra}}
			\end{tabular}
		};

		% ============================================================
		% 15. 右下角落款（画作信息）
		% ============================================================
		\node[anchor=south east, text=qlInk, opacity=0.62, font=\footnotesize\itshape]
			at ([shift={(-1.35,1.25)}]current page.south east)
			{Wang Ximeng, A Thousand Li of Rivers and Mountains, 1113};

		\end{scope}
	\end{tikzpicture}
\end{titlepage}

%
%  配色：ql* 前缀（青绿设色）。本文件用 \providecolor 兜底，
%        单独复制到任何笔记本也能编译；想整体换调改这里即可。
%
%  注意：整幅画面绘于 \begin{scope}[shift={(current page.south west)}] 内，
%        所有坐标以「页面左下角为原点、单位 cm」书写；脱离该 scope
%        直接写 (x,y)，图形会落到页面外而不可见。
%
%  可用字段（在主文件导言区用 \newcommand 预定义即可覆盖缺省值）：
%      \notename     主标题（必填，主模板已定义）
%      \noteauthor   作者（必填，主模板已定义）
%      \notecourseen 眉题英文（缺省 Course Notes）
%      \notesubtitle 副标题（缺省「学习笔记」）
%      \noteextra    底部附加信息（缺省不显示）
% ============================================================

% ---- 《千里江山图》取色（青绿设色） ----
\providecolor{qlSilk}{RGB}{230,218,184}     % 绢底：赭黄
\providecolor{qlSilkDark}{RGB}{208,192,152} % 绢底：暗部
\providecolor{qlGreen2}{RGB}{150,186,128}   % 石绿（淡，远峰）
\providecolor{qlGreen}{RGB}{82,140,106}     % 石绿（主峰）
\providecolor{qlBlue}{RGB}{40,94,132}       % 石青（主峰阴面）
\providecolor{qlOchre}{RGB}{178,134,86}     % 赭石（坡脚/滩渚）
\providecolor{qlInk}{RGB}{56,54,50}         % 墨（点缀/标题）
\providecolor{qlSun}{RGB}{238,214,150}      % 微茫的日色

% ---- 缺省字段（主文件已定义的同名命令不会被覆盖） ----
\providecommand{\notecourseen}{Course Notes}
\providecommand{\notesubtitle}{学\ 习\ 笔\ 记}
\providecommand{\noteextra}{}

\begin{titlepage}
	\begin{tikzpicture}[remember picture, overlay]
		\begin{scope}[shift={(current page.south west)}]

		% ============================================================
		%  1. 绢底（赭黄，下缘略深）
		% ============================================================
		\shade[top color=qlSilk, bottom color=qlSilkDark] (0,0) rectangle (21,29.7);

		% ============================================================
		%  2. 微茫的日色
		% ============================================================
		\fill[qlSun, opacity=0.35] (16.6,26.2) circle[radius=2.6];
		\fill[qlSun, opacity=0.18] (16.6,26.2) circle[radius=4.0];

		% ============================================================
		%  3. 远峰（淡石绿，没入江天）
		% ============================================================
		\fill[qlGreen2, opacity=0.45]
			(0,15.0) -- (2.2,16.2) -- (3.8,15.2) -- (5.6,16.4) -- (7.6,15.0)
			-- (9.8,16.6) -- (12.2,15.2) -- (14.6,16.4) -- (17.0,15.0)
			-- (19.2,16.0) -- (21,14.8) -- (21,12.4) -- (0,12.4) -- cycle;

		% ============================================================
		%  4. 主峰群（石绿 → 石青，自远及近）
		% ============================================================
		% 左峰
		\shade[top color=qlGreen2, bottom color=qlGreen]
			(0,8.4) -- (1.4,10.6) -- (2.6,12.6) -- (3.6,13.6) -- (4.6,12.0)
			-- (5.6,13.0) -- (6.8,11.2) -- (8.0,9.4) -- (9.0,8.2)
			-- (9.4,6.4) -- (9.0,0) -- (0,0) -- cycle;
		% 中峰（最高）
		\shade[top color=qlGreen, bottom color=qlBlue]
			(7.6,8.6) -- (8.8,11.4) -- (9.8,13.8) -- (10.6,15.4) -- (11.4,16.2)
			-- (12.4,14.6) -- (13.4,15.6) -- (14.4,13.2) -- (15.4,11.0)
			-- (16.2,9.0) -- (16.6,7.0) -- (16.4,0) -- (7.6,0) -- cycle;
		% 右峰
		\shade[top color=qlGreen2, bottom color=qlGreen]
			(15.8,8.6) -- (16.8,10.4) -- (17.8,12.2) -- (18.6,13.4) -- (19.6,11.8)
			-- (20.4,12.6) -- (21,10.8) -- (21,0) -- (15.8,0) -- cycle;

		% ============================================================
		%  5. 山体皴线与峰顶苔点
		% ============================================================
		\foreach \x/\y/\l in {2.2/10.4/1.6, 3.6/12.2/1.8, 5.2/11.6/1.4,
		                      9.6/12.6/2.0, 11.0/14.6/2.2, 12.8/13.8/1.8,
		                      14.6/11.6/1.6, 17.6/11.0/1.5, 19.2/11.4/1.3}{
			\draw[qlBlue, opacity=0.30, line width=0.9pt]
				(\x,\y) .. controls (\x+0.22,\y-\l*0.5) and (\x-0.14,\y-\l*0.8) .. (\x+0.12,\y-\l);
		}
		\foreach \x/\y in {3.6/13.7, 5.6/13.1, 9.8/13.9, 11.4/16.3, 13.4/15.7,
		                   18.6/13.5, 20.4/12.7}{
			\fill[qlInk, opacity=0.55] (\x,\y) ellipse[x radius=0.22, y radius=0.12];
		}

		% ============================================================
		%  6. 江水（留白 + 细笔水纹）
		% ============================================================
		\foreach \y in {5.8, 5.0, 4.2, 3.4}{
			\draw[qlBlue, opacity=0.28, line width=0.7pt]
				plot[domain=0.8:20.2, samples=44, smooth]
					(\x, {\y + 0.11*sin(deg(2.4*\x + 4*\y))});
		}

		% ============================================================
		%  7. 滩渚、村落与渔舟
		% ============================================================
		\fill[qlOchre, opacity=0.75]
			(0,0) -- (0,2.4) -- (2.0,3.0) -- (4.2,2.2) -- (6.4,2.8)
			-- (8.6,2.0) -- (10.8,2.6) -- (13.0,1.8) -- (15.4,2.4)
			-- (17.8,1.6) -- (21,2.2) -- (21,0) -- cycle;
		\foreach \x/\y in {2.4/3.0, 3.1/3.0, 2.75/3.5, 13.6/2.5, 14.3/2.5, 13.95/3.0}{
			\draw[qlInk, opacity=0.6, line width=0.8pt] (\x,\y) -- (\x,\y+0.62);
			\fill[qlInk, opacity=0.5] (\x,\y+0.72) circle[radius=0.15];
		}
		\fill[qlInk, opacity=0.85]
			(10.4,4.5) .. controls (11.0,4.28) and (11.9,4.28) .. (12.5,4.5)
			.. controls (12.0,4.74) and (10.9,4.74) .. (10.4,4.5) -- cycle;
		\draw[qlInk, opacity=0.85, line width=0.8pt] (11.45,4.5) -- (11.45,5.2);

		% ============================================================
		%  8. 标题衬底（一层薄绢，保证文字可读）
		% ============================================================
		\fill[qlSilk, opacity=0.55]
			(2.6,16.9) -- (17.4,17.3) -- (18.0,22.6) -- (16.4,23.6) -- (2.2,22.9) -- cycle;

		% ============================================================
		%  9. 细边框（墨线）
		% ============================================================
		\draw[qlInk, line width=0.7pt, opacity=0.45]
			([shift={(0.85,0.85)}]current page.south west) rectangle
			([shift={(-0.85,-0.85)}]current page.north east);

		% ============================================================
		% 10. 顶部眉题（课程英文小字 + 自适应墨线）
		% ============================================================
		\node[anchor=north, text=qlInk, align=center]
			(qlEyebrow) at ([yshift=-3.4cm]current page.north) {%
			{\sffamily\fontsize{12}{15}\selectfont \uppercase{\notecourseen}}
		};
		\draw[qlInk, line width=1.0pt] ($(qlEyebrow.west)+(-0.25cm,0)$) -- ++(-1.1cm,0);
		\draw[qlInk, line width=1.0pt] ($(qlEyebrow.east)+(0.25cm,0)$) -- ++(1.1cm,0);

		% ============================================================
		% 11. 主标题（居中，使用 \notename）
		% ============================================================
		\node[anchor=center, align=center, text=qlInk]
			at ([yshift=-8.3cm]current page.north) {%
			{\fontsize{38}{48}\sffamily\bfseries \notename}
		};

		% ============================================================
		% 12. 菱形 + 双细线
		% ============================================================
		\coordinate (C) at ([yshift=-11.6cm]current page.north);
		\draw[qlInk, line width=1.1pt] (C) ++(-2.2cm,0) -- ++(1.55cm,0);
		\draw[qlInk, line width=1.1pt] (C) ++(0.65cm,0) -- ++(1.55cm,0);
		\node[fill=qlInk, minimum size=0.24cm, inner sep=0pt, rotate=45] at (C) {};

		% ============================================================
		% 13. 副标题（使用 \notesubtitle）
		% ============================================================
		\node[anchor=north, text=qlInk, align=center]
			at ([yshift=-12.4cm]current page.north) {
			{\fontsize{15}{20}\sffamily \notesubtitle}
		};

		% ============================================================
		% 14. 底部作者、日期、附加信息
		% ============================================================
		\coordinate (B) at ([yshift=2.7cm]current page.south);
		\draw[qlInk, line width=0.8pt, opacity=0.8] (B) ++(-1.1cm,0.95cm) -- ++(2.2cm,0);
		\node[anchor=south, text=qlInk, align=center] at (B) {
			\begin{tabular}{c}
				{\fontsize{19}{24}\sffamily\bfseries \noteauthor} \\[0.35cm]
				{\large \sffamily \today}
				\ifstrempty{\noteextra}{}{\\[0.3cm]{\normalsize\sffamily \noteextra}}
			\end{tabular}
		};

		% ============================================================
		% 15. 右下角落款（画作信息）
		% ============================================================
		\node[anchor=south east, text=qlInk, opacity=0.62, font=\footnotesize\itshape]
			at ([shift={(-1.35,1.25)}]current page.south east)
			{Wang Ximeng, A Thousand Li of Rivers and Mountains, 1113};

		\end{scope}
	\end{tikzpicture}
\end{titlepage}
